\chapter{Frobenius 定理：从分布到积分叶} \label{chap:math-geometric-structures}

向量场给出“允许运动的方向”，但给定若干方向并不保证它们真能拼成低维子流形。例如在三维空间中，每点选一个二维平面场，这些平面可能像螺旋一样彼此扭转，因而不存在处处与它们相切的曲面。Frobenius 定理精确回答可积性问题：一个常秩分布局部来自子流形，当且仅当它在 Lie 括号下闭合；用余切语言说，当且仅当它的湮灭一形式生成的微分理想在 $\dd$ 下闭合。

\section{分布、积分流形与 involutivity}

\begin{workedexample}[不可积的平面分布]
在 \(\mathbb R^3\) 上令 \(\alpha=\dd z-x\dd y\)，并令二维分布 \(D=\ker\alpha\)。判断它能否在每一点附近成为某族曲面的切平面。

\textbf{解。} 向量场 \(X=\partial_x\)、\(Y=\partial_y+x\partial_z\) 都满足 \(\alpha(X)=\alpha(Y)=0\)，且处处线性无关，因而张成 \(D\)。由交换子的定义，
\[
[X,Y]f=X(Yf)-Y(Xf)=\partial_z f,\qquad [X,Y]=\partial_z.
\]
但 \(\alpha(\partial_z)=1\)，所以交换子不属于 \(D\)。若 \(D\) 是曲面的切平面，沿曲面切向量场的交换子仍必须切于曲面，矛盾。等价地，\(\dd\alpha=-\dd x\wedge\dd y\) 且 \(\alpha\wedge\dd\alpha=-\dd z\wedge\dd x\wedge\dd y\ne0\)，违背 Frobenius 条件。
\end{workedexample}


秩为 $k$ 的光滑分布是切丛的秩 $k$ 光滑子丛
\[
\mathcal D\subset TM.
\]
局部可以选光滑向量场 $X_1,\ldots,X_k$ 使
\[
\mathcal D_p
=\operatorname{span}\{X_1(p),\ldots,X_k(p)\}.
\]
若存在浸入的 $k$ 维子流形 $L\subset M$，满足每点
\[
T_pL=\mathcal D_p,
\qquad p\in L,
\]
则称 $L$ 为 $\mathcal D$ 的积分流形。若每一点附近都能由这样的积分流形形成局部叶层，则称 $\mathcal D$ 可积。

分布称为 对合的，若任意局部截面 $X,Y\in\Gamma(\mathcal D)$ 都满足
\begin{equation}\label{eq:math-dg-frobenius-involutive}
[X,Y]\in\Gamma(\mathcal D).
\end{equation}
这个条件完全由光滑结构决定，不需要联络或度量。

\begin{theorem}[Frobenius 定理]\label{thm:math-dg-frobenius}
设 $\mathcal D\subset TM$ 是常秩 $k$ 的光滑分布。以下条件在局部等价：
\begin{enumerate}[label=(\roman*)]
\item $\mathcal D$ 可积；
\item $\mathcal D$ 对合的，即满足\eqnrefs{eq:math-dg-frobenius-involutive}；
\item 每点附近存在坐标 $(y^1,\ldots,y^n)$，使
\[
\mathcal D
=\operatorname{span}\{\partial_{y^1},\ldots,\partial_{y^k}\};
\]
\item 若 $\alpha^1,\ldots,\alpha^{n-k}$ 是湮灭丛
\[
\mathcal D^\circ
:=\{\alpha\in T^*M:\alpha|_{\mathcal D}=0\}
\]
的一组局部标架，则存在一形式 $\eta^a{}_b$ 使
\begin{equation}\label{eq:math-dg-frobenius-pfaff}
\dd\alpha^a=\eta^a{}_b\wedge\alpha^b.
\end{equation}
\end{enumerate}
\end{theorem}

从 (iii) 到 (i) 是直接的：固定 $y^{k+1},\ldots,y^n$ 即得到积分叶。从 (i) 到 (ii) 也立即成立，因为切于同一子流形的向量场的 Lie 括号仍切于该子流形。真正困难的是由括号闭合构造适配坐标。

\section{微分理想与 Pfaff 形式}

向量场版本的 involutivity 可以完全转写到余切空间。若 $\alpha\in\Gamma(\mathcal D^\circ)$ 且 $X,Y\in\Gamma(\mathcal D)$，则
\[
\alpha(X)=\alpha(Y)=0.
\]
外微分公式给出
\begin{equation}\label{eq:math-dg-frobenius-dalpha-bracket}
\dd\alpha(X,Y)
=-\alpha([X,Y]).
\end{equation}
因此 $[X,Y]\in\mathcal D$ 当且仅当每个湮灭一形式的 $\dd\alpha$ 在两个 $\mathcal D$ 方向上都为零。在线性代数上，这等价于 $\dd\alpha$ 至少含一个来自 $\mathcal D^\circ$ 的因子，从而得到\eqnrefs{eq:math-dg-frobenius-pfaff}。

若余维为一，$\mathcal D=\ker\alpha$ 且 $\alpha$ 处处非零，Frobenius 条件进一步压缩成
\begin{equation}\label{eq:math-dg-frobenius-codim-one}
\boxed{\alpha\wedge\dd\alpha=0.}
\end{equation}


\par
\begin{workedexample}
在 \(\RR^3\) 中考虑一形式
\[
\alpha=\dd z-x\,\dd y.
\]
求 \(\alpha\wedge\dd\alpha\)，据此判断 \(\ker\alpha\) 是否可积。再自行构造一个非零一形式 \(\widetilde\alpha\) 使 \(\widetilde\alpha\wedge\dd\widetilde\alpha=0\)，并写出其局部积分曲面。
\textbf{解。} \(\dd\alpha=-\dd x\wedge\dd y\)，所以
\[
\alpha\wedge\dd\alpha=(\dd z-x\,\dd y)\wedge(-\dd x\wedge\dd y)
=-\dd z\wedge\dd x\wedge\dd y\ne0,
\]
故 \(\ker\alpha\) 不可积。可取 \(\widetilde\alpha=\dd z\)，它非零且闭，因此 \(\widetilde\alpha\wedge\dd\widetilde\alpha=0\)；积分曲面为 \(z=C\)。
\end{workedexample}
\par

\begin{exbox}[接触型平面场不可积]
在 $\RR^3$ 上取
\[
\alpha=\dd z-x\,\dd y,
\qquad
\mathcal D=\ker\alpha.
\]
有
\[
\dd\alpha=-\dd x\wedge\dd y,
\]
因此
\[
\alpha\wedge\dd\alpha
=-(\dd z-x\dd y)\wedge\dd x\wedge\dd y
=-\dd z\wedge\dd x\wedge\dd y\neq0.
\]
所以 $\mathcal D$ 不可积：不存在处处以这些二维平面为切平面的曲面族。这一例子说明“每点都给出一个二维平面”远不足以保证它们能拼成曲面。
\end{exbox}

余维一 Pfaff 条件的另一端不是“稍微不可积”，而是一个非常刚性的极大不可积结构。设 $M^{2m+1}$ 为奇数维流形，若一形式 $\alpha$ 满足
\begin{equation}
\boxed{
\alpha\wedge(\dd\alpha)^m\neq0
\quad\text{处处成立},
}
\label{eq:math-dg-contact-condition}
\end{equation}
则称 $\alpha$ 为 接触形式，超平面分布
\[
\mathcal H:=\ker\alpha
\]
称为 接触 分布。条件\eqnrefs{eq:math-dg-contact-condition}立刻说明 $\mathcal H$ 不可能满足 Frobenius：若 $\dd\alpha=\eta\wedge\alpha$，则 $(\dd\alpha)^m$ 每一项都含 $\alpha$ 因子，从而 $\alpha\wedge(\dd\alpha)^m=0$，与接触条件矛盾。

更精确地，$\dd\alpha$ 在 $\mathcal H$ 上是非退化二形式。若 $X\in\mathcal H_p$ 且
\[
\dd\alpha_p(X,Y)=0
\qquad\forall Y\in\mathcal H_p,
\]
则由线性代数可推出 $X=0$。因此每个 接触 hyperplane 自身携带一个 $2m$ 维 辛 向量空间 结构
\begin{equation*}\left(\mathcal H_p,\dd\alpha_p|_{\mathcal H_p}\right).\end{equation*}
这说明 接触几何 可以看成“奇数维流形上随点变化的一族 辛 超平面”。

接触形式 还唯一确定一个横向向量场，称为 Reeb 矢量场 $R_\alpha$：
\begin{equation*}\boxed{
\alpha(R_\alpha)=1,
\qquad
\iota_{R_\alpha}\dd\alpha=0.
}\end{equation*}
存在唯一性来自分解
\[
T_pM=\mathcal H_p\oplus\RR R_\alpha(p).
\]
第一条件固定横向分量，第二条件利用 $\dd\alpha|_{\mathcal H}$ 的非退化性消除任何可能的水平修正。

Contact 分布 括号-生成 的充要条件由此给出：经反复 Lie 括号 后张成整个切空间（\eqnrefs{eq:math-dg-contact-step-two-bracket-generating}）。

这个现象与 Frobenius 恰好相反。Frobenius 条件要求 Lie 括号 不产生新方向，于是运动被困在低维叶片中；括号-生成 条件则要求反复 Lie 括号 最终张成整个切空间。Chow--Rashevskii 定理说明，在连通流形上若一个光滑分布 括号 生成，则任意两点都可由分段光滑的水平曲线连接。若再在 $\mathcal H$ 上给定内积，便得到 sub-Riemann 几何：距离只计算水平曲线长度，而“缺失方向”由非交换运动间接产生。

最基本的局部标准型由 接触 Darboux 定理 给出。任一点附近都存在坐标
\[
(q^1,\ldots,q^m,p_1,\ldots,p_m,z)
\]
使 接触形式 在乘以一个处处非零函数后可写成
\begin{equation*}\alpha
=\dd z-\sum_{a=1}^m p_a\,\dd q^a.\end{equation*}
因此 接触几何 与 辛几何 一样没有局部标量曲率型不变量：所有 接触 结构 在足够小邻域内都有相同标准形，真正的信息主要来自全局拓扑、Reeb 动力学 与附加结构。

接触形式 还可以通过增加一个实坐标变成真正的 辛形式。令
\[
\widetilde M:=\RR_s\times M,
\qquad
\lambda:=\ee^s\alpha,
\]
并定义
\begin{equation}
\boxed{
\omega:=\dd\lambda
=\ee^s(\dd s\wedge\alpha+\dd\alpha).
}
\label{eq:math-dg-contact-symplectization}
\end{equation}
由 $\dd^2=0$ 有 $\dd\omega=0$。更重要的是
\begin{align*}
\omega^{m+1}
&=(m+1)\ee^{(m+1)s}
\dd s\wedge\alpha\wedge(\dd\alpha)^m,
\end{align*}
由 接触 条件处处非零。因此 $\omega$ 非退化，$(\RR\times M,\omega)$ 是 辛 流形，称为 $(M,\alpha)$ 的 辛化。接触几何 因而不是与 辛几何 平行的孤立理论，而是奇数维边界数据在增加一个尺度方向后的偶数维 辛 几何。

向量场 $\partial_s$ 是这个 精确 辛 流形 的 Liouville 矢量场，因为
\begin{equation*}\iota_{\partial_s}\omega
=\ee^s\alpha=\lambda,
\qquad
\mathcal L_{\partial_s}\omega=\omega.\end{equation*}
所以沿 $s$ 方向平移会把 辛形式 作指数缩放。另一方面，Reeb 场在 辛化 中恰好由径向 Hamiltonian 识别：若
\[
H(s,x)=\ee^s,
\]
并采用 Hamiltonian 约定
\[
\iota_{X_H}\omega=-\dd H,
\]
则
\begin{equation*}X_H=R_\alpha.\end{equation*}
因为
\[
\iota_{R_\alpha}\omega
=\ee^s\iota_{R_\alpha}(\dd s\wedge\alpha+\dd\alpha)
=-\ee^s\dd s
=-\dd(\ee^s).
\]
这给 Reeb 动力学 一个非常具体的解释：它是在 辛化 上由纯尺度 Hamiltonian 产生的 Hamilton 流。

在 接触 流形 本身，也可为任意光滑函数 $H\in C^\infty(M)$ 定义 接触 Hamiltonian 矢量场 $X_H$。固定约定
\begin{equation}
\boxed{
\alpha(X_H)=H,
\qquad
\iota_{X_H}\dd\alpha
=(R_\alpha H)\alpha-\dd H.
}
\label{eq:math-dg-contact-hamiltonian}
\end{equation}
第二式在 $\mathcal H$ 上唯一确定 $X_H$ 的水平分量，第一式再确定 Reeb 分量。由 Cartan 公式立刻得到
\begin{equation*}\boxed{
\mathcal L_{X_H}\alpha
=(R_\alpha H)\alpha.
}\end{equation*}
因此 接触 Hamiltonian 流 一般不保持具体的一形式 $\alpha$，而保持其 核 所定义的 接触 分布；当 $R_\alpha H=0$ 时才严格保持 $\alpha$。

Hamilton 向量场的 接触 lift 由此确定为 $a=-R_\alpha H$（\eqnrefs{eq:math-dg-contact-hamiltonian-lift}）。

\begin{exbox}[三维 Heisenberg 接触 结构]
在 $\RR^3_{(x,y,z)}$ 上取
\[
\alpha=\dd z-\frac12(x\,\dd y-y\,\dd x).
\]
则
\[
\dd\alpha=-\dd x\wedge\dd y,
\qquad
\alpha\wedge\dd\alpha=-\dd z\wedge\dd x\wedge\dd y\neq0.
\]
一组水平向量场可取
\[
X=\partial_x-\frac y2\partial_z,
\qquad
Y=\partial_y+\frac x2\partial_z,
\]
而 Reeb 场为 $R_\alpha=\partial_z$。直接计算
\[
[X,Y]=\partial_z=R_\alpha.
\]
所以被禁止的竖直方向恰由两个水平运动的交换子生成。这是 Heisenberg 群、非完整约束与 sub-Riemann 控制理论共同的局部模型。
\end{exbox}

\subsection{正则叶片与适配度规}

可积分布给出局部叶片以后，度规还须确定法向距离、切向重标记与积分测度。

\begin{definition}[正则叶状分解]
\label{def:ghs-regular-foliation}
设 $U\subset M$ 是开区域，$I\subset\RR$ 是区间，$\Sigma$ 是 $(d-1)$ 维流形。若存在微分同胚 $F:I\times\Sigma\longrightarrow U$，则 $\Sigma_t=F(\{t\}\times\Sigma)$ 称为一族正则叶。函数 $t:U\to I$ 在每点满足 $\dd t\neq 0$，因此每个 $\Sigma_t$ 都是正则超曲面。
\end{definition}

此定义的用处是将“量子化面”与特定的球面分开。譬如平直空间去掉原点后有 $U\simeq\RR\times S^{d-1}$；而一般紧致流形可以没有全局非临界实函数，因此不存在这样的全局展开。临界点、拓扑变化与边界必须另行处理，不能一律塞进所谓的径向时间。

\begin{proposition}[叶状度规的分解]
\label{prop:ghs-foliation-metric}
在 $I\times\Sigma$ 上任取适配坐标 $(t,y^a)$，正定度规唯一写成
\begin{equation}\label{eq:ghs-foliation-adm}
 g=N^2\dd t^2+h_{ab}(\dd y^a+N^a\dd t)(\dd y^b+N^b\dd t),
\end{equation}
其中 $N>0$ 是 法向进动系数，$N^a$ 是 切向位移，$h_{ab}$ 是 $\Sigma_t$ 上的诱导度规。体积测度为
\begin{equation}\label{eq:ghs-foliation-volume}
 \dd\mu_g=N\sqrt{h}\,\dd t\,\dd^{d-1}y.
\end{equation}
\end{proposition}
\begin{proof}
按 $t$ 和 $y$ 分块写度规矩阵。由于 $h_{ab}$ 正定可逆，令 $N_a=g_{ta}$、$N^a=h^{ab}N_b$，再令 $N^2=g_{tt}-h^{ab}N_aN_b$。正定性和 Schur 补定理给出 $N^2>0$。展开 \cref{eq:ghs-foliation-adm} 得到原矩阵全部分量。分块行列式公式
\[
 \det g=\det h\,(g_{tt}-N_ah^{ab}N_b)=N^2\det h
\]
便证明测度公式。唯一性来自矩阵分块分解的唯一性。
\end{proof}

把 $t$ 增加 $\dd t$ 时，坐标向量 $\partial_t$ 一般不是单位法向。单位法向和法向微分为
\[
 n=N^{-1}(\partial_t-N^a\partial_a),\qquad
 \nabla_n\varphi=N^{-1}(\partial_t-N^a\partial_a)\varphi.
\]
这决定了标量场在两片叶之间真正改变了多少。不能在存在 切向位移 时仅用 $\partial_t\varphi$ 作为正交时间导数。


\section{从局部可积到全局叶空间：holonomy 与主丛水平分布}

Frobenius 定理是局部定理。若 $\mathcal D$ 对合的，通过每一点的局部积分流形可以沿重叠区域唯一延拓，并得到一个极大的连通积分叶；这些极大叶组成 foliation $\mathcal F$。但“每点都有局部叶”并不保证商空间
\[
 M/\mathcal F
\]
本身是流形，甚至不保证 Hausdorff。全局问题取决于叶的闭性、和乐 以及叶在 $M$ 中的回返方式，而这些信息都不包含在点态条件 $[X,Y]\in\Gamma(\mathcal D)$ 中。

最简单的反例是二维环面上的 Kronecker foliation。取
\[
 M=\mathbb T^2=\mathbb R^2/(2\pi\mathbb Z)^2,
 \qquad
 X=\partial_x+\alpha\partial_y,
 \qquad
 \alpha\notin\mathbb Q,
\]
并令 $\mathcal D=\operatorname{span}\{X\}$。秩一分布自动 对合的，因此每条积分曲线
\[
 t\longmapsto
 (x_0+t,\,y_0+\alpha t)\pmod{2\pi}
\]
都是一片叶。由于 $\alpha$ 无理，每片叶在 $\mathbb T^2$ 中稠密。于是叶空间中的单点原像并不是闭集，商空间甚至不是 $T_1$，因而不可能是 Hausdorff 流形。这个例子把“局部 Frobenius 坐标存在”和“存在良好的全局叶坐标”严格区分开。

“叶空间不一定良好”并不意味着任何全局纤维化都不可能。一个重要充分条件由 Ehresmann 纤维化 定理 给出。

\begin{theorem}[Ehresmann fibration theorem]\label{thm:math-dg-ehresmann-fibration}
设 $F:M\to B$ 是光滑、满射、真 淹没。则 $F$ 是局部平凡光滑纤维丛：对每个 $b_0\in B$，存在邻域 $U$ 和微分同胚
\begin{equation}
F^{-1}(U)\cong U\times F^{-1}(b_0)
\label{eq:math-dg-ehresmann-local-trivialization}
\end{equation}
使 $F$ 对应第一坐标投影。
\end{theorem}

淹没 条件保证每个 纤维 都是同维光滑子流形；真性 则提供把局部水平提升延拓到整个小邻域所需的全局控制。取任意 Riemann 度量，并令
\[
\mathcal V=\ker\dd F,
\qquad
\mathcal H=\mathcal V^\perp.
\]
对 $B$ 上的向量场 $Y$，$\dd F|_{\mathcal H}$ 逐点是同构，所以存在唯一水平提升 $Y^H$。若没有 真性，这些水平提升可能在有限时间逃向无穷远，使不同 纤维 无法由流一致识别；真性 恰好排除这一失控。

推论：纤维丛结构。当 $F:M\to B$ 是 真 淹没 且 $B$ 连通时，所有 纤维 微分同胚；结合 long 正合序列 of 同伦群给出纤维丛的同伦分类——Ehresmann 定理是 纤维化 在光滑范畴的版本：真 淹没 自动是 locally trivial 纤维化。例：Hopf 纤维化 $F:S^3\to S^2$（$F(z_1,z_2)=(2\operatorname{Re}(z_1\bar z_2),\,2\operatorname{Im}(z_1\bar z_2),\,|z_1|^2-|z_2|^2)$）中 $S^3$ 紧、$F$ 淹没，自动 真，给出 $S^1$-纤维丛；水平分布对应 Hopf 联络，其曲率是 $S^2$ 上面积形式的拉回，正是 Berry 几何相的数学来源。与叶状结构对比：Ehresmann 的纤维丛结构是'被动'的——从 真 淹没 推出纤维自动形成局部平凡族；叶状结构理论则是'主动'的——从可积分布构造叶，但叶空间可能具有非流形拓扑（如 Reeb 叶状结构中含区间与 Cantor 集混合）。若 $B$ 单连通，纤维丛的整体结构由任一 纤维 的同伦型决定，这把纤维丛分类问题化为同伦论问题，是 特征类 理论的出发点。物理上，量子绝热系统中参数空间 $B$ 上的 Hilbert 丛投影是 淹没，Ehresmann 给出局部平凡化，Berry 联络即水平分布的选择、曲率即 Berry 曲率；参数空间含简并点时 真性 可能失效，对应量子相变。

所以 Frobenius、Ehresmann 与 和乐 处在不同层次：Frobenius 判断一个给定分布是否局部来自叶；Ehresmann 从一个 真 淹没 推出 纤维 全局形成局部平凡族；和乐 则记录沿闭路运输后 纤维 数据如何回返。三者不能互相替代。

主丛联络给出一个更结构化的全局版本。设
\[
 \pi:P\to B
\]
是主 $G$-丛，$\mathfrak g$ 为 Lie 代数。垂直分布是
\[
 \mathcal V=\ker\dd\pi,
\]
而主联络一形式 $\omega\in\Omega^1(P,\mathfrak g)$ 定义水平分布
\begin{equation*}\mathcal H:=\ker\omega,
 \qquad
 TP=\mathcal H\oplus\mathcal V.
\end{equation*}
曲率二形式为
\begin{equation*}\Omega
 =\dd\omega+\frac12[\omega\wedge\omega].
\end{equation*}
若 $X^H,Y^H$ 是水平向量场，则 $\omega(X^H)=\omega(Y^H)=0$，故
\begin{equation}\label{eq:math-dg-curvature-nonintegrability}
 \boxed{
 \Omega(X^H,Y^H)
 =-\omega([X^H,Y^H]).}
\end{equation}
因此
\begin{equation}\label{eq:math-dg-horizontal-frobenius}
 \boxed{
 \mathcal H\text{ involutive}
 \quad\Longleftrightarrow\quad
 \Omega=0.}
\end{equation}
这里的等价使用了主联络曲率本来就是 水平 的事实。曲率不再只是“平行移动绕小回路后的误差”，它就是水平分布违反 Frobenius 条件的精确张量。

固定主丛点 $p\in P$，所有以 $b=\pi(p)$ 为基点的闭路通过水平提升给出结构群元素；这些元素组成 和乐群
\[
\operatorname{Hol}_p(\omega)\subset G.
\]
其中由可缩闭路生成的连通部分称为 限制 和乐，记作
\[
\operatorname{Hol}_p^0(\omega).
\]
平坦性 $\Omega=0$ 只保证 限制 和乐 平凡；若基空间非单连通，完整 和乐 仍可能来自基本群的离散表示。因此局部曲率与全局 单值性 必须分开。

曲率为什么会生成 和乐？考虑基流形上一块由向量 $X,Y$ 张成的无穷小平行四边形，沿其边界作平行移动。若面积尺度为 $\varepsilon^2$，则回到起点后的群元素具有展开
\[
\operatorname{Hol}_{\partial\square_\varepsilon}
=I-\varepsilon^2\Omega(X^H,Y^H)+O(\varepsilon^3),
\]
其中把 Lie 代数 元素通过相应表示解释为无穷小群作用。于是曲率正是 和乐 的二阶局部生成元。

Ambrose--Singer 定理把这一局部图像升级为完整 Lie 代数结论。记 $\mathfrak{hol}_p$ 为 限制 和乐群 的 Lie 代数，则
\begin{equation*}\boxed{
\mathfrak{hol}_p
=\operatorname{span}
\left\{
\operatorname{Ad}_{g_\gamma^{-1}}
\Omega_q(X^H,Y^H)
\right\}
}\end{equation*}
其中 $q$ 遍历可由从 $p$ 出发的水平路径到达的点，$g_\gamma$ 表示沿路径把曲率值平行搬回 $p$ 的群元素。换言之，不是在一个点读取 $\Omega_p$ 就能得到全部 和乐；必须把各点曲率沿路径搬回同一 纤维 后再取 Lie 代数 张成。

对 Levi--Civita 联络，这一结论把 Riemann 曲率与切空间中的 和乐群 联系起来。若 和乐 严格小于 $SO(n)$，就意味着存在被平行移动保持的额外张量。例如 Kähler 几何中的复结构 $J$ 满足 $\nabla J=0$，因此 和乐 降到 $U(m)$；Calabi--Yau 情形进一步降到 $SU(m)$。这里不展开特殊 和乐 分类，但逻辑已经清楚：
\[
\text{曲率}
\Longrightarrow
\text{holonomy Lie algebra}
\Longrightarrow
\text{平行张量与额外几何结构}.
\]

Frobenius 定理因此应当分成两个层次理解：局部层次由 Lie 括号或微分理想决定，保证适配坐标和积分叶存在；全局层次还要检查叶空间拓扑与 和乐。下一章把几何结构允许显含时间，研究张量、度量与由度量决定的联络在演化中的精确变化律。
闭合于 Lie 括号是方向分布能够拼成积分叶的关键条件。下面把这个条件与非可积反例对照，检查局部坐标能否沿允许方向铺开。

\section{哪些方向能够拼成曲面}

在 \(\RR^3\) 每一点选一个二维平面，单独看每个平面都没有问题。然而若想找曲面，使这些平面恰好是曲面的切平面，就必须检查沿两个允许方向先后走动会不会产生第三个方向。

\begin{definition}[光滑分布与积分流形]
设 \(M\) 为 \(n\) 维光滑流形。秩为 \(r\) 的光滑分布 \(\mathcal D\) 给每点 \(p\in M\) 指定 \(r\) 维子空间 \(\mathcal D_p\subset T_pM\)，且在每点附近可由 \(r\) 个光滑向量场逐点张成。若 \(r\) 维浸入子流形 \(N\) 满足 \(T_pN=\mathcal D_p\)，则称 \(N\) 为积分流形。这里浸入指局部参数化 \(F:U\subset\RR^r\to M\) 的微分在每点单射。
\end{definition}

例如 \(\operatorname{span}\{\partial_x,\partial_y\}\) 的积分曲面是 \(z=\text{常数}\)。反例是
\[
X=\partial_x,\qquad Y=\partial_y+x\partial_z,\qquad
\mathcal D=\operatorname{span}\{X,Y\}.
\]
二者处处独立，\([X,Y]=\partial_z\) 却不在分布内：比较 \(\partial_x,\partial_y\) 系数，若 \(\partial_z=aX+bY\)，必有 \(a=b=0\)，与 \(\partial_z\) 系数矛盾。

这个交换子还能直接量出“小闭合路”未能闭合的方向。从
\((x_0,y_0,z_0)\) 出发，依次沿 \(X\) 走 \(s\)、沿 \(Y\) 走 \(t\)、沿
\(-X\) 走 \(s\)、沿 \(-Y\) 走 \(t\)。前两步使 \(x\) 变为 \(x_0+s\)，
所以沿 \(Y\) 的 \(z\) 增量为 \((x_0+s)t\)；后两步中沿
\(-Y\) 的增量为 \(-x_0t\)。最后 \(x,y\) 回到原值而
\(z\) 多了 \(st\)。若允许方向真是同一曲面的切方向，这条由曲面内
曲线拼成的闭合坐标路不可能跑出曲面；非零的 \(st\partial_z\)
正是交换子的最低阶效应。换句话说，若两个场的流总可交换，
这样拼出的路径必回到起点；净位移非零正说明流不可交换，
其大小在最低阶由 Lie 括号量出。

\begin{definition}[对合分布]
若每一对取值于 \(\mathcal D\) 的局部光滑向量场 \(U,V\) 都满足 \([U,V](p)\in\mathcal D_p\)，称 \(\mathcal D\) 对 Lie 括号\emph{对合}。只检查一组局部基即可，因为
\([fU,gV]=fg[U,V]+f(Ug)V-g(Vf)U\)。
\end{definition}

\begin{theorem}[Frobenius 局部可积定理]
设 \(\mathcal D\) 是光滑流形 \(M\) 上的常秩光滑分布。对每个 \(p\in M\)，存在经过 \(p\) 的局部 \(r\) 维积分流形，当且仅当 \(\mathcal D\) 对合。等价地，对合时在 \(p\) 附近存在局部坐标 \((u^1,\ldots,u^r,v^1,\ldots,v^{n-r})\)，使
\(\mathcal D=\operatorname{span}\{\partial_{u^1},\ldots,\partial_{u^r}\}\)。
\end{theorem}

\begin{proof}
先证必要性。若 \(N\) 是积分流形，取局部方程 \(v^1=\cdots=v^{n-r}=0\) 表示 \(N\)。任意切于 \(N\) 的场 \(U,V\) 满足 \(Uv^a=Vv^a=0\) 在 \(N\) 上成立。对限制在 \(N\) 上恒零的光滑函数作切向导数仍为零，故 \([U,V]v^a=U(Vv^a)-V(Uv^a)=0\)。因此交换子也切于 \(N\)。

再证充分性。围绕 \(p\) 取坐标 \((x^1,\ldots,x^r,y^1,\ldots,y^{n-r})\)，使 \(\mathcal D_p\) 到前 \(r\) 个坐标方向的投影可逆。缩小坐标域后，可把分布的基唯一写成
\[
X_i=\partial_{x^i}+\sum_{a=1}^{n-r}A_i^a(x,y)\partial_{y^a},
\qquad 1\le i\le r.
\]
直接展开 \([X_i,X_j]\)，其 \(\partial_x\) 分量为零。对合性又要求该括号属于 \(\mathcal D\)；而 \(\mathcal D\) 中非零向量必有非零 \(\partial_x\) 分量，所以 \([X_i,X_j]=0\)。

令 \(\Phi_i^t\) 为 \(X_i\) 的局部流。这里由交换子为零推出流可交换，理由可以具体写出：沿 \(X_i\) 的流对 \(X_j\) 作推前，其对流参数的导数是 \([X_i,X_j]\) 的推前，因此 \((\Phi_i^s)_*X_j=X_j\)。于是两条曲线 \(t\mapsto\Phi_i^s\Phi_j^t(q)\) 与 \(t\mapsto\Phi_j^t\Phi_i^s(q)\) 都满足同一条 \(X_j\) 的常微分方程，且在 \(t=0\) 有相同初值；局部解唯一性给 \(\Phi_i^s\Phi_j^t=\Phi_j^t\Phi_i^s\)。固定横截面上的初值 \((0,y_0)\)，定义
\[
F(u,y_0)=\Phi_1^{u^1}\circ\cdots\circ\Phi_r^{u^r}(0,y_0).
\]
交换性保证 \(\partial_{u^i}F=X_i(F)\)。\(F\) 在原点的微分把 \(u\) 方向送到 \(X_i(p)\)，把 \(y_0\) 方向送到 \(\partial_{y^a}\)；这 \(n\) 个向量线性无关。由反函数定理，\((u,y_0)\) 是局部坐标。固定 \(y_0\) 后，\(u\mapsto F(u,y_0)\) 就是所求积分流形。
\end{proof}

证明中的关键是对合性使局部基变成可交换的流，因此固定横截面初值时，先走哪一个方向都不改变所得曲面。

局部坐标结论还给出一种实际寻找曲面的方法。把上述 \(v^a\) 看成 \(n-r\) 个函数，则 \(Xv^a=0\) 对每个取值于 \(\mathcal D\) 的向量场 \(X\) 成立，而 \(\dd v^1,\ldots,\dd v^{n-r}\) 线性无关。这样的函数叫作分布的\emph{局部第一积分}：沿允许方向走动时，它们的值不变。反过来，若找到了这组函数，映射 \(v=(v^1,\ldots,v^{n-r})\) 的微分满秩，由隐函数定理，每个非空水平集 \(v=c\) 都是 \(r\) 维嵌入子流形，且切空间为
\(\ker \dd v^1\cap\cdots\cap\ker \dd v^{n-r}=\mathcal D\)。这证明对合性、适配坐标与足够多的局部第一积分是同一件事的三种表达。它并不承诺 \(v^a\) 能在整个流形上选成单值函数。

“浸入”与“嵌入”也不能混为一谈。浸入只要求局部参数化的微分单射；嵌入还要求参数化到像的拓扑同胚，使像在附近确实像一片不自交的坐标曲面。Frobenius 定理给的是每点附近的嵌入叶，若把局部叶沿允许方向尽量延长，得到的全局叶可能仅是浸入的。取二维环面 \(\mathbb T^2=\RR^2/(2\pi\ZZ)^2\)，令
\(X=\partial_x+\alpha\partial_y\)，其中 \(\alpha\) 为无理数。秩一分布 \(\operatorname{span}\{X\}\) 自动对合，其积分曲线为
\(t\mapsto(x_0+t,y_0+\alpha t)\bmod 2\pi\)。若它在不同 \(t,s\) 处重合，则 \(t-s=2\pi m\)、\(\alpha(t-s)=2\pi n\)，无理性迫使 \(m=n=0\)，所以曲线没有自交；但它在环面中稠密。稠密性是无理旋转的标准结论：固定目标后取 \(t=x-x_0+2\pi m\)，第二坐标随整数 \(m\) 走过无理旋转 \(2\pi\alpha m\) 的轨道；把圆等分后由抽屉原理可在轨道中找到任意短的非零步长，从而逼近圆上任意点。稠密性的关键是无理旋转轨道在圆周上稠密：将圆分成足够多等弧，抽屉原理给出一个非零轨道步长模 \(2\pi\) 任意小；这个小步长的整数倍可在圆周上逼近任意目标点，而这些点仍属于同一轨道闭包。因此任意目标 \((x,y)\) 都可由某列轨道点逼近。作为整个像，这条曲线并非嵌入子流形。局部可积不等于能用一个全局方程把叶整齐分层。

余维为一时常用处处非零的一形式 \(\alpha\) 写 \(\mathcal D=\ker\alpha\)。对于 \(U,V\in\Gamma(\mathcal D)\)，外微分定义给
\(\dd\alpha(U,V)=-\alpha([U,V])\)。因此 \(\mathcal D\) 对合当且仅当 \(\dd\alpha\) 在 \(\mathcal D\) 上为零。取 \(Z\) 满足 \(\alpha(Z)=1\)，并取 \(U,V\in\ker\alpha\)，则
\((\alpha\wedge \dd\alpha)(Z,U,V)=\dd\alpha(U,V)\)。在每点，交替三形式确实由这一类取值决定：把任意三向量各自分解为 \(Z\) 方向分量与 \(\ker\alpha\) 分量，\(\alpha\) 只读到 \(Z\) 分量，含两个 \(Z\) 分量的交叉项成对抵消，剩下的只含 \(\ker\alpha\) 分量的取值。故条件又等价于 \(\alpha\wedge \dd\alpha=0\)。取 \(\alpha=\dd z-x\,\dd y\)，则 \(\dd\alpha=-\dd x\wedge \dd y\)，
\(\alpha\wedge \dd\alpha=-\dd z\wedge \dd x\wedge \dd y\ne0\)。它的核正是上面的 \(X,Y\) 所张成的分布，因而不存在以它为切平面的局部曲面。

把约束改为 \(\beta=\dd z-\dd(xy)=\dd z-y\,\dd x-x\,\dd y\)，结论就变了：
\(\beta=\dd(z-xy)\)，所以 \(\dd\beta=0\)，Frobenius 条件成立；
固定 \(z-xy=c\) 的曲面就是积分流形。其切向量可取
\(\partial_x+y\partial_z\) 与 \(\partial_y+x\partial_z\)，
代入 \(\beta\) 都为零。这一正反对照也提醒我们：
\(\dd\beta=0\) 是充分条件，但对余维一分布并非必要条件，
必要且充分的是较弱的 \(\beta\wedge \dd\beta=0\)。



\par
\begin{workedexample}
先复算 \([X,Y]=\partial_z\)；再把约束改为 \(\alpha=\dd z-\dd(xy)\)，验证 \(\alpha\wedge \dd\alpha=0\)，并直接写出积分曲面。最后说明在变秩分布上为何不能原样套用上述定理。（可选）把环面无理流的稠密性用抽屉原理补全：把圆分成 \(N\) 段，说明轨道点中必有两点的模一距离小于 \(1/N\)。
\textbf{解.} 前面的 \(X=\partial_x\)、\(Y=\partial_y+x\partial_z\) 满足 \([X,Y]=\partial_z\)，而 \(\partial_z\notin\operatorname{span}\{X,Y\}\)，故原分布不对合。改取 \(\alpha=\dd z-\dd(xy)\)，则 \(\dd\alpha=0\)，所以 \(\alpha\wedge\dd\alpha=0\)。其核由约束 \(z-xy=C\) 的切向量张成，故水平集是积分曲面。变秩分布的秩会改变，无法组成定秩子丛，Frobenius 定理的假设失效。可选环面部分：\(N+1\) 个轨道点落入 \(N\) 段，抽屉原理给出两点模一距离不超过 \(1/N\)；两点之差产生任意小的非零转角，其整数倍可逼近任意圆周点，故无理转动的轨道稠密。
\end{workedexample}
\par


\begin{theorem}[Frobenius 局部积分叶定理]
设 $\mathcal D\subset TM$ 为光滑常秩 $k$ 分布，局部截面在 Lie 括号下封闭。每点附近存在坐标使\[\mathcal D=\operatorname{span}\{\partial_{y^1},\ldots,\partial_{y^k}\}.\]固定其余坐标得到局部积分叶。
\end{theorem}

\begin{proof}
目标是在分布对合的邻域上构造适配坐标使 $\mathcal D$ 由前 $k$ 个坐标方向张成。对秩 $k$ 归纳：$k=1$ 时，取非零局部截面 $X_1$，流盒定理给出坐标使
\[
X_1=\partial_{y^1},
\]
故积分曲线就是一维叶。

设结论对秩 $k-1$ 成立。取 $p$ 附近非零截面 $X_1\in\Gamma(\mathcal D)$，再次用流盒坐标令
\[
X_1=\partial_{y^1}.
\]
令横截面
\[
\Sigma:=\{y^1=0\}.
\]
通过对其余局部生成元减去适当的 $X_1$ 分量，可以选取
\[
Y_2,\ldots,Y_k\in\Gamma(\mathcal D)
\]
使它们在 $\Sigma$ 上切于 $\Sigma$。于是
\[
\mathcal E:=\mathcal D\cap T\Sigma
\]
是 $\Sigma$ 上秩 $k-1$ 分布。若 $Y,Z$ 是 $\mathcal E$ 的截面，则它们既切于 $\Sigma$ 又属于 $\mathcal D$。切于固定子流形的向量场在 Lie 括号下封闭，而 $\mathcal D$ 假设对合，因此
\[
[Y,Z]\in\Gamma(\mathcal D\cap T\Sigma)=\Gamma(\mathcal E).
\]
由归纳假设，缩小邻域后可在 $\Sigma$ 上取适配坐标
\[
(y^2,\ldots,y^k,y^{k+1},\ldots,y^n),
\]
使 $\mathcal E$ 由前 $k-1$ 个坐标方向张成；固定
$y^{k+1},\ldots,y^n$ 得到 $\mathcal E$ 的局部积分叶族 $L_c$。特别地，其中一片 $L_0$ 通过 $p$。

现在令 $\Phi_t$ 为 $X_1$ 的局部流。关键是证明它保持分布 $\mathcal D$。取局部标架 $X_1,\ldots,X_k$ 张成 $\mathcal D$。对合条件给出
\[
[X_1,X_a]=B^b{}_aX_b
\]
其中 $B^b{}_a$ 为光滑函数。沿 $X_1$ 的积分曲线，推前标架满足线性常微分方程；等价地，$(\Phi_{-t})_*X_a$ 的导数仍是这些向量的线性组合。因此初值位于 $\mathcal D$ 的向量在整个局部流中始终留在 $\mathcal D$，即
\[
(\Phi_t)_*\mathcal D=\mathcal D.
\]

对 $\Sigma$ 中每一片局部叶 $L_c$ 定义
\[
L_c^{\mathrm{sat}}
:=\{\Phi_t(q):q\in L_c,\ |t|<\varepsilon\}.
\]
其切空间由 $X_1$ 与 $(\Phi_t)_*T_qL_c$ 张成；前者属于 $\mathcal D$，后者由于分布被流保持也属于 $\mathcal D$。维数同为 $k$，故
\[
T_{\Phi_t(q)}L_c^{\mathrm{sat}}
=\mathcal D_{\Phi_t(q)}.
\]
因此这些饱和叶给出 $p$ 附近的局部叶层。把流参数作为 $y^1$，保留 $\Sigma$ 上的
$y^2,\ldots,y^n$，便得到局部坐标，其中
\[
\mathcal D
=\operatorname{span}\{\partial_{y^1},\ldots,\partial_{y^k}\}.
\]
这就得到 (iii)，并同时构造了局部积分叶。

$\mathbb R^3$ 中的分布 $\mathcal D=\operatorname{span}\{\partial_x,-y\partial_z+\partial_y\}$：$[\partial_x,-y\partial_z+\partial_y]=0\in\mathcal D$，对合，积分叶为 $z+\tfrac12y^2=\text{const}$ 的抛物柱面。

反例 $\mathcal D=\operatorname{span}\{\partial_x,x\partial_z+\partial_y\}$：$[\partial_x,x\partial_z+\partial_y]=\partial_z\notin\mathcal D$，不对合，不存在通过每点的二维积分曲面——这是接触几何的典型例子。

Frobenius 定理在控制论中给出可达性条件：向量场张成的分布对合当且仅当系统可沿积分叶任意移动。
\end{proof}

\begin{proposition}[余维一的 Pfaff 可积判据]
设 $\alpha$ 为处处非零光滑一形式。以下条件局部等价：$\ker\alpha$ 对合、$\alpha\wedge\dd\alpha=0$、存在 $\eta$ 使 $\dd\alpha=\eta\wedge\alpha$，以及存在非零 $\mu$ 和局部坐标 $u$ 使 $\alpha=\mu\dd u$。
\end{proposition}

\begin{proof}
余维一时分布 $\mathcal D=\ker\alpha$ 由单一一形式 $\alpha$ 确定。Frobenius 条件 $\dd\alpha=\eta\wedge\alpha$ 等价于 $\alpha\wedge\dd\alpha=0$，后者是可直接检验的代数条件。关键观察是：在 $n$ 维流形上，$\alpha\wedge\dd\alpha$ 是一个三形式，三形式的纤维维数为 $\binom n3$（仅 $n=3$ 时为一维），因此其消失条件只含 $\binom n3$ 个分量函数，远少于检验全部 Lie 括号。

先证 Frobenius $\Rightarrow$ Pfaff。若 $\dd\alpha=\eta\wedge\alpha$，则
\begin{equation*}
 \alpha\wedge\dd\alpha
 =\alpha\wedge\eta\wedge\alpha=0,
\end{equation*}
因为 $\alpha\wedge\alpha=0$（一形式与自身的楔积为零）。这一方向不需要 $\eta$ 的任何具体性质，只要它存在即可。

再证 Pfaff $\Rightarrow$ Frobenius。假设 $\alpha\wedge\dd\alpha=0$。局部取向量场 $N$ 使 $\alpha(N)=1$（归一化），并定义一形式
\begin{equation*}
 \eta:=-\iota_{N}\dd\alpha,
\end{equation*}
即 $\eta$ 是 $\dd\alpha$ 沿 $N$ 方向的缩并（带负号）。这里 $N$ 的存在性用到 $\alpha$ 处处非零，从而可在每点的余切空间中选一个对偶向量；局部光滑性由 $\alpha$ 光滑保证。

需在基向量对上逐一验证 $\dd\alpha=\eta\wedge\alpha$：取 $X,Y\in\ker\alpha$，并把 $\{X,Y,N\}$ 当作局部标架（在余维一情形下 $\ker\alpha$ 处处补上 $N$ 即给出 $TM$）：
\begin{itemize}
\item 在 $(X,Y)$ 上：由 $\alpha\wedge\dd\alpha=0$ 在 $(N,X,Y)$ 上取值，展开
\begin{equation*}
(\alpha\wedge\dd\alpha)(N,X,Y)
=\alpha(N)\dd\alpha(X,Y)-\alpha(X)\dd\alpha(N,Y)+\alpha(Y)\dd\alpha(N,X),
\end{equation*}
后两项因 $\alpha(X)=\alpha(Y)=0$ 消失，只剩 $\alpha(N)\dd\alpha(X,Y)=\dd\alpha(X,Y)=0$。而 $(\eta\wedge\alpha)(X,Y)=\eta(X)\alpha(Y)-\eta(Y)\alpha(X)=0$。两式相等。
\item 在 $(X,N)$ 上：$(\eta\wedge\alpha)(X,N)=\eta(X)\alpha(N)-\eta(N)\alpha(X)=\eta(X)=-\dd\alpha(N,X)=\dd\alpha(X,N)$。两式相等。这里最后一步用了 $\eta$ 的定义与二形式的反对称性。
\item 在 $(N,N)$ 上：两二形式在重复向量上取值均为零。
\end{itemize}
在所有基向量对上取值相同，故
\begin{equation*}
 \dd\alpha=\eta\wedge\alpha.
\end{equation*}

在余维一可积情形下，Frobenius 定理给出局部坐标 $(x^1,\ldots,x^{n-1},u)$ 使
\begin{equation*}
\alpha=\mu(x,u)\,\dd u,
\end{equation*}
其中 $\mu$ 是处处非零的光滑函数。叶状结构由 $u=\text{const}$ 给出，每片叶是 $n-1$ 维子流形。等价地，存在局部积分因子 $\mu^{-1}$ 使 $\mu^{-1}\alpha$ 成为恰当形式 $\dd u$；这正是经典 Pfaff 方程 $\alpha=0$ 求解的第一积分存在性。

接触结构的反例：在 $\RR^3$ 上取 $\alpha=\dd z-x\,\dd y$，则 $\dd\alpha=-\dd x\wedge\dd y$，于是
\begin{equation*}
\alpha\wedge\dd\alpha
=-(\dd z-x\,\dd y)\wedge\dd x\wedge\dd y
=-\dd z\wedge\dd x\wedge\dd y\neq0.
\end{equation*}
Pfaff 判据立即给出 $\ker\alpha$ 不可积，不存在二维积分曲面。这正是接触几何的出发点：$\alpha\wedge\dd\alpha\neq0$ 处处成立时 $\alpha$ 称为接触形式，$\ker\alpha$ 称为接触分布，它"最大限度地不可积"。

叶状结构的正例：在 $\RR^3$ 上取 $\alpha=\dd z-x\,\dd x$，则 $\dd\alpha=-\dd x\wedge\dd x=0$，从而 $\alpha\wedge\dd\alpha=0$。由于 $\alpha=\dd(z-\tfrac12x^2)$ 是恰当形式，等值面 $z-\tfrac12x^2=c$ 的切空间正是 $\ker\alpha$，因而它们构成所求的抛物柱面叶状结构。

于是\eqnrefs{eq:math-dg-frobenius-codim-one}与 Frobenius 条件局部等价。Pfaff 判据 $\alpha\wedge\dd\alpha=0$ 的优势在于不需要显式构造 $\eta$，只需检查一个三形式的消失；在余维一情形下它把可积性问题完全化为一个代数条件，这是接触几何与叶状结构理论分歧的精确分水岭。
\end{proof}

\begin{proposition}[接触分布的两步括号生成性]
设 $(M^{2m+1},\alpha)$ 为接触流形，$m\ge1$，$\mathcal H=\ker\alpha$。每点满足\[\mathcal H_p+[\mathcal H,\mathcal H]_p=T_pM,\qquad\alpha([X,Y])=-\dd\alpha(X,Y).\]
\end{proposition}

\begin{proof}
对水平向量场 $X,Y\in\Gamma(\mathcal H)$，有 $\alpha(X)=\alpha(Y)=0$。由\eqnrefs{eq:math-dg-frobenius-dalpha-bracket}，
\begin{equation}
\alpha([X,Y])=-\dd\alpha(X,Y).
\label{eq:math-dg-contact-bracket-vertical}
\end{equation}
固定一点 $p$ 和任意非零 $X_p\in\mathcal H_p$。由于接触条件等价于 $\dd\alpha|_{\mathcal H_p}$ 非退化，作为 $\mathcal H_p$ 上的辛形式，对每个非零 $X_p$ 必存在 $Y_p\in\mathcal H_p$ 使
\begin{equation*}
\dd\alpha_p(X_p,Y_p)\neq0.
\end{equation*}
具体地，取 $\mathcal H_p$ 的一组辛基 $\{e_i,f_j\}$ 使 $\dd\alpha_p(e_i,f_j)=\delta_{ij}$，将 $X_p$ 展开为 $X_p=\sum_i(a_i e_i+b_i f_i)$，则选 $Y_p=\sum_i(-b_i e_i+a_i f_i)$ 即给出
\begin{equation*}
\dd\alpha_p(X_p,Y_p)=\sum_i(a_i^2+b_i^2)=\|X_p\|^2>0.
\end{equation*}
把 $X_p,Y_p$ 延拓为局部水平场 $X,Y$（即 $\alpha(X)=\alpha(Y)=0$ 在邻域内成立），便由\eqnrefs{eq:math-dg-contact-bracket-vertical}得到
\begin{equation*}
\alpha_p([X,Y]_p)=-\dd\alpha_p(X_p,Y_p)\neq0.
\end{equation*}

设 $R_\alpha$ 为 Reeb 向量场，满足 $\alpha(R_\alpha)=1$、$\iota_{R_\alpha}\dd\alpha=0$。由 $\alpha_p([X,Y]_p)\neq0$，可写
\begin{equation*}
[X,Y]_p
=\alpha_p([X,Y]_p)\,R_\alpha(p)+H_p,
\qquad H_p\in\mathcal H_p,
\end{equation*}
所以 $[X,Y]_p$ 具有非零 Reeb 方向分量。又 $\mathcal H_p=\ker\alpha_p$ 维数为 $2m$，而 $T_pM$ 维数为 $2m+1$，故
\begin{equation*}
\dim\bigl(\mathcal H_p+[\mathcal H,\mathcal H]_p\bigr)
\ge\dim\mathcal H_p+1
=2m+1
=\dim T_pM,
\end{equation*}
反向不等式自动成立，即
\begin{equation}
\boxed{
\mathcal H_p+[\mathcal H,\mathcal H]_p=T_pM.
}
\label{eq:math-dg-contact-step-two-bracket-generating}
\end{equation}
接触分布因而是两步括号-生成分布：虽然单个允许速度永远限制在 $\mathcal H$ 中，但通过"先沿 $X$、再沿 $Y$、再反向返回"的小回路，可在二阶产生横向位移。
\end{proof}

\begin{proposition}[接触 Hamilton 场的辛化提升]
采用 $\alpha(X_H)=H$、$\iota_{X_H}\dd\alpha=(R_\alpha H)\alpha-\dd H$。在辛化 $\omega=\dd(\ee^s\alpha)$ 上按 $\iota_{X_{\widetilde H}}\omega=-\dd\widetilde H$ 定义场，则 $\widetilde H=\ee^sH$ 满足\[X_{\widetilde H}=X_H-(R_\alpha H)\partial_s.\]
\end{proposition}

\begin{proof}
在辛化上令
\[
\widetilde H(s,x)=\ee^sH(x).
\]
把 Hamiltonian 矢量场分解为
\[
X_{\widetilde H}=X_H+a\,\partial_s.
\]
利用\eqnrefs{eq:math-dg-contact-symplectization}，逐项计算
\[
\iota_{X_{\widetilde H}}\omega
=\ee^s\left[
a\alpha-H\dd s+\iota_{X_H}\dd\alpha
\right].
\]
而
\[
-\dd\widetilde H
=-\ee^s(H\dd s+\dd H).
\]
比较两式并代入\eqnrefs{eq:math-dg-contact-hamiltonian}，得到
\[
a=-R_\alpha H.
\]
所以
\begin{equation}
\boxed{
X_{\widetilde H}
=X_H-(R_\alpha H)\partial_s.
}
\label{eq:math-dg-contact-hamiltonian-lift}
\end{equation}

对常数 $\lambda>0$，记 $r_\lambda(s,x)=(s+\ln\lambda,x)$ 为沿 $s$ 方向的平移。由 $\widetilde H$ 的定义，
\begin{equation*}
\widetilde H\circ r_\lambda
=\ee^{s+\ln\lambda}H(x)
=\lambda\,\widetilde H(s,x),
\end{equation*}
即 $\widetilde H$ 是沿 $s$ 方向的 $1$ 次齐次函数。

对应地，辛 形式 $\omega=\dd(\ee^s\alpha)$ 在 $r_\lambda$ 下满足 $r_\lambda^*\omega=\lambda\,\omega$，因此 Hamiltonian 流保持这一齐次结构：若 $\phi_t$ 是 $\widetilde H$ 的流，则 $r_\lambda\circ\phi_t\circ r_{\lambda^{-1}}$ 也是 $\widetilde H$ 的流，由唯一性得 $r_\lambda\circ\phi_t=\phi_t\circ r_\lambda$。

Reeb 流作为特例：取 $H\equiv1$，则 $R_\alpha H=0$ 且 $X_H=R_\alpha$（由接触 Hamiltonian 的定义方程 $\iota_{X_H}\dd\alpha=(R_\alpha H)\alpha-\dd H=0$ 与 $\alpha(X_H)=H=1$）。此时
\begin{equation*}
X_{\widetilde H}=R_\alpha,
\end{equation*}
即 Reeb 流恰好提升为 辛化 上与 $s$ 坐标无关的 Hamiltonian 流。这给出 Reeb 流所有接触不变量的 辛 解释。

三维接触情形的显式公式：在标准接触 $\RR^3$ 上取 $\alpha=\dd z-\tfrac12(x\,\dd y-y\,\dd x)$，Reeb 场 $R_\alpha=\partial_z$。对 $H=H(x,y,z)$，接触 Hamiltonian 场为
\begin{equation*}
X_H
=\left(\frac{\partial H}{\partial y}+\frac{x}{2}\frac{\partial H}{\partial z}\right)\partial_x
-\left(\frac{\partial H}{\partial x}-\frac{y}{2}\frac{\partial H}{\partial z}\right)\partial_y
+\left(H-\frac{x}{2}\frac{\partial H}{\partial x}-\frac{y}{2}\frac{\partial H}{\partial y}\right)\partial_z,
\end{equation*}
提升到 $\RR^4_{(s,x,y,z)}$ 后沿 $s$ 方向的分量为 $-\partial_z H$。当 $H$ 与 $z$ 无关时提升为齐次 Hamiltonian 系统，这正是经典 Hamilton 力学在接触框架下的对应。

接触 Hamilton 系统也能编码摩擦，但必须保持同一符号约定。
在 $\alpha=\dd z-p\dd q$ 上，本章采用 $\alpha(X_H)=H$，因此
\[
\dot q=-H_p,\qquad\dot p=H_q+pH_z,
\qquad\dot z=H-pH_p,
\qquad\dot H=(R_\alpha H)H.
\]
若要得到 $\ddot q+\gamma\dot q+\omega^2q=0$、$\gamma>0$，应取
\[
H=-\left(\tfrac12p^2+\tfrac12\omega^2q^2+\gamma z\right).
\]
于是 $\dot q=p$、$\dot p=-\omega^2q-\gamma p$。
机械能 $E=\tfrac12p^2+\tfrac12\omega^2q^2$ 满足
\[
\dot E=-\gamma p^2\le0,\qquad\dot H=-\gamma H.
\]
若使用原来的正号 $H=E+\gamma z$，在当前约定下得到的是反摩擦，
不能把 $\dot H=\gamma H$ 解释成摩擦能量耗散。
辛化后 $\widetilde H=\ee^sH$ 保守，且 $\dot s=-R_\alpha H=\gamma$；
额外坐标抵消接触能量的指数因子，保留 $\widetilde H$。

接触 Hamiltonian 动力学 因而可以看成一类对 $s$ 方向齐次的 辛 Hamiltonian 动力学；Reeb 流 对应 $H\equiv1$ 的特例，而一般情形下 Reeb 分量 $-R_\alpha H$ 正好度量系统偏离保守的程度。
\end{proof}

\begin{theorem}[Ehresmann 局部平凡化定理]
设 $F:M\to B$ 为光滑、满射、真的下沉映射。每个 $b_0$ 附近存在邻域 $U$ 与微分同胚\[\Phi:U\times F^{-1}(b_0)\to F^{-1}(U),\qquad F\circ\Phi(b,x)=b.\]
\end{theorem}

\begin{proof}
在 $b_0$ 附近选坐标球 $U$ 使其闭包紧，并对每个 $b\in U$ 取从 $b_0$ 指向 $b$ 的径向路径 $\gamma_b:[0,1]\to U$，满足 $\gamma_b(0)=b_0$、$\gamma_b(1)=b$。把 $\dot\gamma_b(t)$ 水平提升到 $M$，从任意 $x\in F^{-1}(b_0)$ 出发解 ODE
\[
\dot{\widetilde\gamma}_{b,x}(t)
=\dot\gamma_b(t)^H,
\qquad
\widetilde\gamma_{b,x}(0)=x.
\]
沿提升曲线计算 $F$ 的变化率，由链式法则与水平提升的定义，
\begin{align*}
\frac{\dd}{\dd t}F(\widetilde\gamma_{b,x}(t))
&=\dd F_{\widetilde\gamma_{b,x}(t)}\bigl(\dot{\widetilde\gamma}_{b,x}(t)\bigr)\\
&=\dd F_{\widetilde\gamma_{b,x}(t)}\bigl(\dot\gamma_b(t)^H\bigr)
=\dot\gamma_b(t),
\end{align*}
最后一步用到 $\dd F|_{\mathcal H}$ 是恒等提升。积分两端并代入初值 $F(\widetilde\gamma_{b,x}(0))=F(x)=b_0=\gamma_b(0)$，得到
\[
F(\widetilde\gamma_{b,x}(t))=\gamma_b(t),
\]
因此提升曲线始终停留在对应纤维之上。因为 $F$ 真，紧小球 $\overline U$ 的原像 $F^{-1}(\overline U)$ 紧，水平 ODE 的解不会在有限时间逃向无穷远；缩小 $U$ 后，所有这些 ODE 在 $t\in[0,1]$ 上统一存在。定义
\[
\Phi(b,x):=\widetilde\gamma_{b,x}(1).
\]
由 ODE 对初值 $x$ 与参数 $b$ 的光滑依赖，$\Phi$ 光滑；且由上式 $F(\Phi(b,x))=\gamma_b(1)=b$。构造逆映射时，对 $(y,b)\in F^{-1}(U)$ 取从 $b$ 沿径向返回 $b_0$ 的路径 $\gamma_b(1-t)$，从 $y$ 出发水平提升至 $t=1$，得到 $F^{-1}(b_0)$ 中的一点，与 $b$ 组合即给出 $\Phi^{-1}$，它同样由 ODE 光滑依赖保证光滑。因此
\[
\Phi:U\times F^{-1}(b_0)\to F^{-1}(U)
\]
是微分同胚，并满足 $F\circ\Phi(b,x)=b$，这正得到\eqnrefs{eq:math-dg-ehresmann-local-trivialization}。
\end{proof}

\begin{proposition}[主联络曲率与水平可积性]
对主联络 $\omega$，令 $\Omega=\dd\omega+\frac12[\omega\wedge\omega]$、$\mathcal H=\ker\omega$。则\[\Omega(X^H,Y^H)=-\omega([X^H,Y^H]),\qquad\Omega=0\Longleftrightarrow\mathcal H\text{ 对合}.\]全局水平截面还要求和乐平凡。
\end{proposition}

\begin{proof}
对任意 $\mathfrak g$-值一形式 $\omega$ 与任意向量场 $X,Y$，外微分的 不变 公式为
\[
 \dd\omega(X,Y)
 =X\{\omega(Y)\}-Y\{\omega(X)\}-\omega([X,Y]).
\]
令 $X=X^H$、$Y=Y^H$ 为水平提升。由水平分布定义 $\mathcal H=\ker\omega$，有 $\omega(X^H)=\omega(Y^H)=0$，前两项均为零，剩下
\[
 \dd\omega(X^H,Y^H)=-\omega([X^H,Y^H]).
\]
再算楔括号项。由定义
\[
 [\omega\wedge\omega](X^H,Y^H)
 =[\omega(X^H),\omega(Y^H)]-[\omega(Y^H),\omega(X^H)],
\]
每个因子均为零，故
\[
 [\omega\wedge\omega](X^H,Y^H)=0.
\]
把两式代入曲率定义 $\Omega=\dd\omega+\tfrac12[\omega\wedge\omega]$，得到
\begin{align*}
\Omega(X^H,Y^H)
&=\dd\omega(X^H,Y^H)+\tfrac12[\omega\wedge\omega](X^H,Y^H)\\
&=-\omega([X^H,Y^H])+0,
\end{align*}
即\eqnrefs{eq:math-dg-curvature-nonintegrability}。

若 $\Omega=0$，则对一切水平对有 $\omega([X^H,Y^H])=0$，即 $[X^H,Y^H]\in\ker\omega=\mathcal H$，括号仍水平，Frobenius 定理给出局部水平叶。

反之若水平分布 对合的，则 $[X^H,Y^H]\in\mathcal H$，故 $\omega([X^H,Y^H])=0$，由上式曲率在全部水平对上消失；主联络曲率又满足 horizontality，即它被水平向量场完全决定，因此 $\Omega=0$ 处处成立，即\eqnrefs{eq:math-dg-horizontal-frobenius}。

第一 Bianchi 恒等式 $D\Omega=0$ 可视为水平分布的"高阶可积性"约束：若 $\Omega=0$ 给出局部叶，则 $D\Omega=0$ 自动满足；若 $\Omega\neq0$，Bianchi 限制曲率沿水平方向变化的可能方式。在联络语言中 Bianchi 恒等式是 $D\Omega=0$。
协变外微分一般并不满足 $D^2=0$：在伴随值形式上
$D^2\eta=[\Omega\wedge\eta]$，其平方恰由曲率控制。
黎曼第一 Bianchi 还使用无挠条件，第二 Bianchi 则是协变曲率恒等式。

当结构群 $G$ 为 Abelian（如 $U(1)$ 电磁联络），$[\omega\wedge\omega]=0$，曲率简化为 $\Omega=\dd\omega$。

此时 Frobenius 障碍就是 $\dd\omega=0$，即 $\omega$ 闭；局部叶存在当且仅当 $\omega$ 局部恰当，这正是电磁势的规范不变性 $\omega\mapsto\omega+\dd f$ 不改变曲率。

非 Abelian 情形下 $[\omega\wedge\omega]$ 项给出非交换性贡献，使曲率不再是线性算子，Yang--Mills 理论的非线性性正源于此。

对 $U(1)$ 联络在圆 $S^1$ 上的平行移动，沿闭路 $\gamma$ 的和乐为
\begin{equation*}
\operatorname{Hol}(\gamma)=\exp\left(i\oint_\gamma\omega\right)
=\exp\left(i\int_D\Omega\right),
\end{equation*}
最后一步要求闭路在相应平凡化内确实为某曲面 $D$ 的边界，才可用 Stokes 定理；圆周上的一般非可缩闭路没有这样的 $D$。曲率 $\Omega$ 因此是和乐的"密度"：把闭路包围的曲率积分即得和乐角度。非 Abelian 情形下此式变为非交换 Wilson 圈
\begin{equation*}
W(\gamma)=\mathcal P\exp\left(\oint_\gamma\omega\right),
\end{equation*}
其中 $\mathcal P$ 表示路径序，反映非交换性使局部贡献不能简单交换。

但 $\Omega=0$ 仍只是局部可积性。固定 $p_0\in P$ 后，沿基流形路径作水平提升定义平行移动。若基空间存在非平凡闭路，平坦联络仍可能具有非平凡和乐表示
\[
 \operatorname{Hol}:\pi_1(B,b_0)\to G.
\]
只有当这些闭路的和乐全部平凡时，从 $p_0$ 出发的水平叶才会在每个基点上选出唯一一点，从而成为全局水平截面。若 $B$ 单连通，则 $\pi_1(B)=0$，平坦性使平行移动与路径无关，于是该局部 Frobenius 结构可以沿整个连通基底拼成全局水平截面。由此可见：曲率控制局部可积性，和乐 控制全局拼接；前者是微分条件，后者是拓扑条件，两者共同决定联络的可积性结构。

曲率与和乐的分离对应特征类的两层结构：第一 Chern 类 $c_1\in H^2(B,\ZZ)$ 通过 de Rham 表示由曲率积分给出（局部信息），但其整体挠部分由 和乐 决定（全局信息）。Chern--Weil 理论把曲率不变多项式映射到上同调类，给出 $c_k$、$p_k$、Euler 类等的微分代表；积分特征类的挠部分不能由实曲率形式检出，需整体转移函数或微分特征类等额外数据。$\hat A$-亏格和通常的数值指标同样不能检出一般挠类，不应把它们当作完整恢复挠信息的工具。
\end{proof}

\begin{proposition}[无穷小矩形和乐的曲率系数]
局部取 $D=\dd+A$，以列向量满足 $\dot U=-A(\dot\gamma)U$。沿正向 $X,Y,-X,-Y$ 的坐标小矩形，\[U_{\partial\square}=I-\varepsilon^2F(X,Y)+O(\varepsilon^3),\qquad F=\dd A+A\wedge A.\]
\end{proposition}

\begin{proof}
取局部联络形式 $A=A_i\,\dd x^i$，沿 $X$、$Y$ 方向各走长度 $\varepsilon$ 的小矩形。平行移动在每条边上由 $\dot U=-A(\dot\gamma)U$ 给出。
为完整保留二阶项，所有 $A$ 及其导数都在起点计算，四条边分别为
\[
\begin{aligned}
U_1&=I-\varepsilon A_X+\tfrac{\varepsilon^2}{2}(A_X^2-\partial_XA_X),\\
U_2&=I-\varepsilon A_Y+\varepsilon^2(\tfrac12A_Y^2-\partial_XA_Y-\tfrac12\partial_YA_Y),\\
U_3&=I+\varepsilon A_X+\varepsilon^2(\tfrac12A_X^2+\tfrac12\partial_XA_X+\partial_YA_X),\\
U_4&=I+\varepsilon A_Y+\tfrac{\varepsilon^2}{2}(A_Y^2+\partial_YA_Y).
\end{aligned}
\]
每式余项均为 $O(\varepsilon^3)$；对列向量，完整乘积为
$U_{\partial\square}=U_4U_3U_2U_1$。
原来只保留每边的一阶项再声称精确到二阶，会遗留错误的 $A_X^2,A_Y^2$ 项；
这些同边项必须与上面的二阶指数项同时计算。
一次项全部消失，二次项整理为
\[
U_{\partial\square}
=I-\varepsilon^2
\left(
\partial_XA_Y-
\partial_YA_X+[A_X,A_Y]
\right)
+O(\varepsilon^3).
\]
括号正是曲率分量
\[
F(X,Y)
=(\dd A+A\wedge A)(X,Y).
\]
因此曲率不是平行移动误差的比喻，而是无穷小闭路和乐的首个非平凡系数。

把四条边的乘积按 $\varepsilon$ 阶数展开。零阶项为 $I$。一阶项为
\begin{equation*}
-\varepsilon A_X-\varepsilon A_Y+\varepsilon A_X+\varepsilon A_Y=0,
\end{equation*}
正反向贡献相消，反映沿闭路回到起点的最低阶必然平凡。二阶项分三类：
\begin{itemize}
\item 同向导数项：$-\varepsilon^2\partial_XA_Y$ 与 $+\varepsilon^2\partial_YA_X$，来自沿一边走时联络值在另一边方向的变化；
\item 交换项：$-\varepsilon^2[A_X,A_Y]$，来自非交换 Lie 代数的乘积顺序差异；
\item 同边往返项：如 $(-\varepsilon A_X)(\varepsilon A_X)=-\varepsilon^2A_X^2$，但这些项在同一条边往返时已被路径序局部解吸收，不出现在最终公式中。
\end{itemize}
三类贡献合并即得曲率分量。

当 $G$ 为 Abelian（如 $U(1)$），$[A_X,A_Y]=0$，曲率简化为 $F_{XY}=\partial_XA_Y-\partial_YA_X$，这正是电磁场强 $F_{\mu\nu}=\partial_\mu A_\nu-\partial_\nu A_\mu$。沿小回路的和乐为
\begin{equation*}
U_{\partial\square}=\exp\left(-i\varepsilon^2F_{XY}\right)+O(\varepsilon^3),
\end{equation*}
对应回路包围的磁通；只有进一步考虑随时间变化的回路或场，才联系到 Faraday 感应定律。

对 $SU(N)$ 联络，交换项 $[A_X,A_Y]$ 给出非线性自相互作用，是 Yang--Mills 方程
\begin{equation*}
D^\mu F_{\mu\nu}=0
\end{equation*}
非线性的来源。物理上这对应胶子自相互作用：与电磁场光子不带电荷不同，胶子本身带色荷，因此曲率含 $A\wedge A$ 项。

把无穷小回路公式积分到有限回路，得到和乐沿任意闭路 $\gamma$ 的表达式
\begin{equation*}
\operatorname{Hol}(\gamma)=\mathcal P\exp\left(-\oint_\gamma A\right),
\end{equation*}
对边界为 $\gamma$ 的曲面 $\Sigma$ 应用非交换 Stokes 定理，得
\begin{equation*}
\operatorname{Hol}(\gamma)=\mathcal P_\Sigma\exp\left(-\int_\Sigma \widehat F\right),\qquad \widehat F(x)=U_x^{-1}F(x)U_x,
\end{equation*}
这里先选择从同一基点到 $x$ 的扫描路径，以 $U_x$ 把各点曲率运输回同一纤维，再采用与扫描相容的曲面积序 $\mathcal P_\Sigma$。

不能直接把不同纤维中的 $F(x)$ 相乘。

把 $\Sigma$ 分成 $N^2$ 个面积为 $1/N^2$ 的小方格并取 $N\to\infty$，每个小方格贡献 $I-\frac1{N^2}F+O(N^{-3})$，乘积收敛到上式。

该式适用于边界为 $\gamma$ 的这类曲面与选定扫描。

Ambrose--Singer 定理则说受限和乐的 Lie 代数由平行运输回基点的曲率值生成；非可缩闭路上的整体和乐仍需独立拓扑数据。

在量子力学中，参数空间上的绝热演化给出 $U(1)$ 联络（Berry 联络）$\mathcal A=i\langle\psi|\dd\psi\rangle$，曲率 $\mathcal F=\dd\mathcal A$ 是 Berry 曲率。沿参数空间闭路演化一周后波函数获得相位 $\exp(i\oint\mathcal A)$，这正是无穷小回路公式在 $U(1)$ 情形的应用。Aharonov--Bohm 效应中电子绕磁通量管的相位差则是非平凡 和乐 的典型例子：联络局部平坦（磁场为零区域）但全局 和乐 非平凡，体现曲率与 和乐 的全局-局部分离。
\end{proof}

